\chapter*{General Conclusion}
\phantomsection
\addcontentsline{toc}{chapter}{General Conclusion}
\markboth{General Conclusion}{General Conclusion}
\label{chap:general-conclusion}

The general conclusion closes the report by returning to the initial problem, objective, methodology, and validated results. It should be short, balanced, and evidence-based.

Begin by restating the problem and the project objective in one paragraph. Then summarize the main work performed, the most important results, and the degree to which the objectives were achieved.

Mention limitations honestly and connect them to future work. Future work should be realistic: propose extensions, deeper validation, deployment improvements, broader experiments, or usability enhancements that follow from the report.

Avoid introducing new technical material, new citations, new results, or emotional closing statements. The conclusion should leave the reader with a clear understanding of what was achieved and what remains to be done.